\documentclass[11pt]{article}
\usepackage{preprint}

\title{Political Economic Analysis on Inequality in America}
\author{Jiadi He\email{jh2288@georgetown.edu}}
\shorttitle{Political Economy of Inequality in America}
\keywords{political economy; power; inequality}

\begin{document}
\maketitle
\keywordsline

The core argument of this paper is to emphasize the impact of political and economic power on American economic inequality, which includes two interrelated factors: one is the political polarization brought by the rise of Republican power. The second is the economic driving force behind the rise of the Republican Party, that is, the return of big capital and the political influence of capital interest groups.

\section{Interaction between political polarization and economic inequality}

The close relationship between American economic inequality and political polarization provides an important clue to understanding the political and economic driving force of inequality.

In the past century, economic inequality and political polarization in the United States have shown a highly positive correlation. The development trend of the two can be said to follow suit and closely follow each other. In the era of “great compression” in the mid-20th century, the degree of political polarization of both houses of Congress was also the lowest. During this period, both parties were controlled by moderates, and the differences between the two parties were rarely kidnapped by ideology. Therefore, cross-voting between Party members of the two parties was a common phenomenon, and the failure of government function caused by ideological dispute was extremely rare. Before and after that, when the economic inequality was serious, it was also the time when the partisan disputes within parliament became white-hot. The most serious division before the 1940s was in the 1920s when the income share of the highest 1\% rose rapidly and peaked before the economic crisis in 1929. Since the late 1970s, the two groups of data have risen rapidly at the same time, exceeding the level of a century ago at the beginning of the 21st century.

So what does the simultaneous rise of economic polarization and political polarization mean since the 1980s? On the one hand, economic inequality directly contributes to political polarization. Naturally, high-income people will invest more time and resources in supporting political parties against redistribution. Of course, the people at the bottom will support another faction. On the other hand, political polarization also promotes conservative Republicans to formulate public policies that increase unequal distribution to obtain political support. If the Republican Party has a majority in Congress, it will use this advantage to formulate policies to reduce redistribution, such as the tax cut reform during Reagan and George W. Bush. If the Republican Party does not have an advantage, it will become a force hindering changes to the status quo, such as tax reform policies to raise taxes for the rich or policies to raise the minimum wage. In a word, in a divided political environment, the Republican Party can effectively use positive legislation and negative resistance to achieve its policy intention, and it is very successful. In this political environment, inequality caused by non-political reasons - such as changes in technology, education, and social habits - cannot be changed through the political process, but strengthened.

But the next question is, if the Republican Party becomes increasingly conservative and extreme, which hinders the formulation of public policy of unequal distribution, why can't the Democratic Party compete? After all, in the institutional environment of bipartisan confrontation such as the United States, if the forces of both sides are balanced, political polarization may lead to the stagnation of policy innovation, but it does not necessarily lead to the right shift of the overall policy. Moreover, before the late 1970s, the democratic party representing the left had ruled American politics for more than 40 years. The policy of protecting equality has penetrated all aspects of American society. What made the Democratic Party collapse so easily?

Since the late 1970s, American politics has not simply become increasingly polarized and left-right confrontation, but polarization in the process of overall deviation to the right. It can be seen that the Republican Party shifted sharply to the right during this period. In contrast, the overall left shift of the Democratic Party is much smaller, and it is mainly concentrated in the south, and the Democrats in the North hardly have much shift. Since the late 1970s, the Democratic Party's camp in the South has lost rapidly. Today, almost all the southern states have become the sphere of influence of conservatives, and most of the northern regions are controlled by the Democratic Party, which makes people think of the North-South Division before the civil war. Moderate Democrats are generally replaced by Republicans, leaving only those Democrats who are tough enough to survive in this increasingly hostile environment.

Another important factor is the change of the Democratic Party and the Liberals themselves. While the Republicans and conservatives rapidly moved away from the redistribution policy, the Democrats and liberals also deviated from the welfare policy they were keen on in the same period but shifted their attention to other more personalized social issues, such as race, environmental protection, gender, sexual orientation, abortion, women's rights and so on. The attention to these social problems has greatly reduced the investment of Democrats and liberals in issues such as income equality and social welfare.

From a larger perspective, the transformation of Democrats and liberals is also the product of changes in American social culture and political culture. Although the equal rights movement and anti-war movement in the 1960s brought unprecedented social equality, they also caused new divisions in American society and politics. In response to the promotion of race and gender equality by Democrats and liberals, the American conservative movement began to recover (i.e.\ Neoconservatism). As early as 1964, marked by the participation and defeat of Republican candidate Barry Goldwater, a large number of senior Republican politicians and Goldwater lost the election at the same time, which also made the young generation of American conservatives rise in the party and replace the organic fraction in the Communist Party in a radical way.

Since then, the Democratic Party has changed from a class-based party to a racial alliance, and the Republican Party has also defined itself as an anti-civil right and anti-immigration party, laying the seeds for the rise of right-wing populism in the future. In 1972, the MC Government Fraser Commission revised the Democratic Party's rules for selecting nominees, requiring some seats for women, ethnic minorities, and young people, which to some extent stimulated the radical forces in the Republican conservatives.

A direct consequence of these changes is that cultural differences have replaced the economic welfare policy and become a new focus of the debate between the two parties. They have also become an element of the growing geographical and ethnic differentiation of voters of the two parties: on the one hand, the north-south differentiation, on the other hand, the differentiation of the white middle class and the rest of the population. For Republicans, cultural differences have become a very effective tool to shift goals when fighting back against Democrats and American society's demand for changing economic inequality policies. For example, the large-scale and conservative evangelicals are an important voter base of the Republican Party. Among them, low-income people are far more likely to vote for Republicans representing high-income groups than nonevangelical believers at the same income level. When assisting bush in his re-election in 2004, Karl Rove, known as “Bush's brain”, successfully incited many states to take anti-homosexual marriage as an amendment to the state constitution, thus stimulating conservative voters to vote.

In short, after the 1970s, the political polarization brought by the changes in American social culture has a far-reaching impact on American economic and welfare policies. Despite the growing differences between the two parties, class differences are no longer the focus of differences. On the contrary, identity politics caused by cultural differences began to dominate American politics. Economic equality and social welfare are still important elements of the Democratic Party's program, but the Democratic Party has been unable to resist the overall right shift of economic policy.

\section{The return of big capital and the rise of capital interest groups}

Although the overall right shift environment of American economic policy has the party and socio-cultural background, it is not a simple cultural and policy issue. In the final analysis, it is a political result, that is, the result of power distribution.

In the 1950s and early 1960s, one of the most important debates in American Social Sciences was who ruled American Society - ordinary citizens, economic elites, or interest groups? Scholars represented by Wright Mills believe that American politics is ruled by a power elite composed of the military, chaebols, and senior government officials. Their decisions determine the structure of American society and affect the lives of middle and lower-class people. But the result of the debate was the victory of Robert Dahl's theory of pluralistic democracy: power is widely distributed among relatively equal and diversified civil groups, interest groups, and political parties, and there is no single absolute group or alliance. These organizations compete with each other to influence the government, to ensure that the basic interests and needs of ordinary citizens can be expressed in public policy.

However, when Dahl put forward his theory in the early 1960s, the United States was in a very special historical stage: the “great compression” era known for economic equality and political compromise he described was very different from the American society before and after. For those Americans who experienced the so-called “roaring twenties”, the surging social movements we see today are just the result of another economic division and political decline in the past century.

Similar to the situation in the early 20th century, the social changes in the early 21st century are dominated by the control of big capital over politics and public policies. Similarly, today's research once again proves that American politics is increasingly controlled by elites and interest groups. The traditional theory of voter control and pluralistic democracy does not exist. In other words, the influence of ordinary voters in American politics almost does not exist. Even among interest groups, there is great inequality. The influence of capital interest groups is twice that of general public organizations. Moreover, interest groups rarely represent the opinions of ordinary voters. Sometimes, for example, on the issue of gun possession, the opinions of the American Gun association are different from those of ordinary citizens, but they can still influence government decision-making.

When analyzing the growing economic inequality, American scholars Jacob Hacker and Paul Pierson called the most important feature of this process “winner-take-all politics” - American politics has become organized combat. In this context, capital interest groups make a comeback and change the distribution and redistribution structure of market income by influencing public policies and government regulations such as financial market, enterprise regulation, labor relations, and tax law, to tilt it to large capital (such as 1\% of high-income people). By the mid-1980s when Reagan was in power, the Republican Party had successfully become “the most powerful political party organization in American history”. To a large extent, this is inseparable from the financial help of big capital. By the mid-1980s, the Republican Party had five times as much money as the Democratic Party.

However, unlike in the early 20th century, today's American capital interest groups no longer rely on the personal influence of a few business giants such as Morgan and Rockefeller but gather under lobbying groups such as the American Chamber of Commerce and the Business Council to express their interests in a larger, more complex and more covert way. In terms of capital, interest groups (such as trade unions) representing the interests of labor and the lower class have never been rivals of capital interest groups. However, this gap has widened in recent years. A 2014 academic study showed that at the federal government level, 84\% of lobbying spending came from large companies and business associations representing large capital.

But money is not the only tool used by capital interest groups. The direct influence of lobbying on politicians and the decision-making process has also become too straightforward to be accepted by today's public. On the contrary, through academic circles and think tanks endorsing their interests and propositions, their propositions and demands can be better expressed by using the “revolving door” formed among academic circles, political circles, and business circles. Their social and academic backgrounds are highly consistent, and their ideology and economic interests are inseparable.

In addition, it has become the most important tactic of today's capital interest groups to mobilize the whole people with the help of various social groups and enhance their public influence. For example, the Koch brothers, a business giant known for its active participation in politics, has funded and even directly established a large number of conservative foundations, think tanks (including the famous Cato Institute and traditional foundation), and civil society organizations, which have a great impact on issues such as tax reduction and climate change. According to a report in the Washington Post, the Americans for Prosperity, a civil society group recently organized by them, can rival the Democrats and Republicans in terms of manpower and capital and can be called the third-largest party in the United States.

For another example, interest groups representing big capitalists will hide their investment identity, organize various grass-roots organizations, and more secretly express their commercial interests under the cloak of these grass-roots organizations. This strategy is ironically called “astroturfing”, but it is extremely effective and has successfully helped them avoid the supervision of the American public.

In addition to using various ways to express their views and safeguard their interests, today's American capital interest groups pay more attention to influencing the decision-making and operation of the US government in terms of ideology, ideology, and even social values, to shift the government to the right in ideology and political culture, and more actively accept and even actively promote the policies advocated by capital interest groups, This makes them more diverse, flexible and hidden in the expression and maintenance of interests, so they are also more effective. This strategy is called “thought laundering”, and its impact and sustainability are far beyond the imagination of business giants in the early 20th century.

In terms of ideology and political culture, one of the most important achievements of American capital interest groups in the past 30 years is to cultivate a people who are quite hostile to the government. After decades of efforts by interest groups, President Reagan's famous words - “the most frightening words in English are: I come from the government; I come to help” and “government expansion means the shrinkage of freedom” - have been deeply rooted in the political subconsciousness of today's American people. Since Reagan, words such as tax cuts, deregulation, and small government have gradually become the Political Correctness in American political culture and a powerful weapon for capital interest groups to control Republican politicians, helping them successfully exclude moderate Republican members from the mainstream camp. In the era of “great compression”, moderates dominated American politics, and “conservative Democrats” and “liberal Republicans” were praising words at that time. Today, these people are regarded as wall riders. In particular, Republicans who agree with the views of the left on some issues will be called “RINO, Republican in name only”, which means nominally Republicans, and will be blacklisted and denounced.

The result of this process is not only to push the Republican Party to extremes. More importantly, under the dual pressure of money and ideology, Democratic politicians are unable to compete in this process. Money plays an increasingly important role in elections, so the cost of leaning to the left is obvious; however, it is unrealistic to rely on the strength of civil society to compete with large capital and capital interest groups - the left-wing social forces have been lack of motivation for a long time since the new deal to form new forces to unite all parties. “Occupy Wall Street movement” can be said to be the first real progressive movement of left-wing social forces since the 1930s, but it is still in its infancy and lacks strong echo at the American elite level, which is far from enough to compete with conservative forces.

Most importantly, succumbing to big capital and capital interest groups is not a partisan issue. The Democratic Party does not have the opportunity to change the status quo and counter the policies of the Republican Party of the United States. During the Clinton period (1992-2000), the Democratic Party controlled both houses of Parliament, so it can take this opportunity to change the status quo. On the contrary, the Clinton administration is not inferior to the Republican administration in promoting neoliberal policies. When Gore participated in the election in 2000 with inequality as the election topic, he was ruthlessly attacked by the party, which can show the right ideological tide in the Democratic Party.

In this context, although the differences between the Republican Party and the Democratic Party are getting deeper and deeper, the substantive differences have become more and more blurred. Because no matter which party enters the White House, the fundamental interests of Wall Street and big capital groups are difficult to touch. More importantly, both parties have completely surrendered ideologically to Wall Street. As mentioned earlier, the natural result is that the income of the richest 1\% did not decline after the financial crisis, but quickly picked up and further increased. Because there is no public policy, the negative impact of the crisis on wealth concentration is short and weak. The tax increase reform during Obama's administration has only made a fine adjustment to the income of the rich without substantial harm. Although Obama tried his best to save wall street from the financial stock disaster, it does not mean that Wall Street bought his account. Shortly after the financial crisis, during a visit to Wall Street, Obama tried to persuade Wall Street to agree to control the rising salaries of managers and help calm public anger. But no Wall Street representative attended his speech. For another example, according to a news in April 2016, Hillary Clinton refused to publish her speech when she met with Goldman Sachs and other Wall Street representatives, which was questioned by Bernie Sanders. The Democratic Party has been so, while the Republican Party has become the agent of this interest group and has been completely captured.

The above analysis gives sufficient evidence to explain the reasons for the surging anti-elite sentiment in American Society: wealth is concentrated to the rich, and money is increasingly linked to power, allowing the rich to use wealth to control politics, replacing the “people” in the ownership, rule by the people and enjoyment of the people with “chaebol”, to return American society to the “plutocracy” at the beginning of the 21st century.

\begin{references}[Reference]
\item Hacker, Jacob S., and Paul Pierson. 2006. \textit{Off Center: The Republican Revolution and the Erosion of American Democracy}. Yale University Press.
\item Hacker, Jacob S., and Paul Pierson. 2010. \textit{Winner-Take-All Politics: How Washington Made the Rich Richer-and Turned Its Back on the Middle Class}. Simon \& Schuster.
\item Perlstein, Rick. 2009. \textit{Before the Storm: Barry Goldwater and the Unmaking of the American Consensus}. Nation Books.
\item Feldmann, Linda, ed. 2019. \textit{Review of Trump’s Unraveling Republican Party: How We Reached This Point}. The Washington Post.
\item Sullivan, Andrew, ed. 2016. \textit{Review of Democracies End When They Are Too Democratic}. New York Magzine.
\item Wright Mills, C. 1956. \textit{The Power Elite}. New York, NY: Oxford University Press.
\item Dahl, Robert A. 2006. \textit{A Preface to Democratic Theory}. Chicago, IL: University of Chicago Press.
\item Dahl, Robert A. 2005. \textit{Who Governs?: Democracy and Power in an American City, Second Edition.} New Haven, CT: Yale University Press.
\item Gilens, Martin, and Benjamin I. Page. 2014. “Testing Theories of American Politics: Elites, Interest Groups, and Average Citizens.” \textit{Perspectives on politics} 12(3): 564–81. \url{http://dx.doi.org/10.1017/s1537592714001595}.
\item Bartels, Larry M. 2018. \textit{Unequal Democracy: The Political Economy of the New Gilded Age - Second Edition}. Princeton, NJ: Princeton University Press.
\item Kwak, James. 2014. “Cultural Capture and the Financial Crisis.” \textit{In Preventing Regulatory Capture}, eds. Daniel Carpenter and David A. Moss. Cambridge: Cambridge University Press, 71–98.
\item Bump, Philip, ed. 2014. \textit{Americans for Prosperity May Be America’s Third-Biggest Political Party}. The Washington Post.
\item Fukuyama, Francis. \textit{Left Out}. The American Interest, 22-28.
\item Meyerson, Harold, ed. 2016. \textit{The Long March of Bernie’s Army}. American Prospect.
\item Bruinius, Harry, ed. 2016. \textit{Why Hillary Holds So Little Allure for Many Millennials?} The Christian Science Monitor.
\end{references}

\end{document}
